\chapter{Experiments}
\label{chap:experiments}
\setcounter{section}{0}

\section{Quantitative Evaluation}

All experiments were conducted on the same machine: an OMEN Gaming Laptop 16 with an AMD Ryzen AI 7 350 @ 2 GHz, an NVIDIA RTX 5070 Laptop GPU (7.71 GB VRAM), 32 GB RAM, running Windows 11 Home (25H2) and always plugged in using official charger. Although as it will be seen, this setup over compensates for the need to run these models, using quantization techniques. Cloud-based systems were evaluated over the same Wi-Fi connection to ensure consistent network conditions.


\subsection{Autopilot System Evaluation}

\subsubsection{Model Selection}

Autopilot relies on two different families of models: Vision-Language Models (VLMs), used by the Scout, Vision, and QA agents, and Large Language Models (LLMs), used by the Intention Agent and the Manager Agent. Each family was benchmarked independently, as they present different objectives in the system.

\paragraph{VLM Selection.}
The three VLM roles in Autopilot (Scout, Vision, and QA) share the same underlying model, but they impose different objectives. Scout and QA tasks can be fulfilled using almost any VLM, as most of them are reasonably good with OCR (Optical Character recognition). However, the Vision Agent presents a different challenge, it must return precise screen coordinates. This visual grounding features, is not frequent in VLMs and therefore it was used for the primary selection of the visual model.

Six candidate models were evaluated across three task templates (\textit{Write a Note}, \textit{Search and Open File}, and \textit{Change Text Format in Word}) [ref with each specific prompt or to the bigger table], which all templates require to identify icons and text in screen. The porcentage result for each template, presents with the overall accuracy. Although more models were tested, the INT 4 versions where the most competitive, and therefore the ones presented in the table. The full list of evaluated models is provided in Appendix A.


\begin{table}[H]
\centering
\begin{tabular}{|l|c|c|c|c|c|c|c|}
\hline
\textbf{Model} & \textbf{Params} & \textbf{Quant.} & \textbf{Write} & \textbf{Search} & \textbf{Format} & \textbf{Inference} & \textbf{Avg Vram} \\
 & & & \textbf{(2 steps)} & \textbf{(5 steps)} & \textbf{(5 steps)} & \textbf{Time (s)} & \textbf{(\%)} \\
\hline
Qwen3 2B VL & 2B & INT4 & 100 & 100 & 100 & 34.71 & 32.96 \\
\hline
Qwen3 4B VL & 4B & INT4 & 100 & 100 & 100 & 38.99 & 48.69 \\
\hline
Mistral 3B  & 3B & INT4 & 0   & 0   & 0   & 33.40 & 47.2  \\
\hline
Gemma3 4B   & 4B & INT4 & 0   & 0   & 0   & 31.07 & 51.89 \\
\hline
\end{tabular}
\caption{VLM candidate evaluation across three task templates.}
\label{tab:vlm_candidates}
\end{table}

Both Qwen3 2B VL and Qwen3 4B VL achieved perfect grounding across all three templates, correctly identifying both text and icon targets in every step. Mistral 3B and Gemma3 4B failed entirely at coordinate prediction, hallucinated positions, and were therefore discarded, despite the inference time and the small VRAM used.

Between the two viable candidates, Qwen3 2B VL and Qwen3 4B VL are identical in accuracy (100\% across all templates) but differ in the inference time and the resources used. Since both models achieve the same grounding performance, the smaller model is preferable under the project's consumer hardware constraint. Qwen3 2B VL INT4 was selected as the VLM for all three agent roles (Scout, Vision, and QA).


\paragraph{LLM Selection.}
The LLM must perform two distinct roles: intent classification with the Intention Agent, and task orchestration with the Manager Agent. The model must produce the correct intent name and all required parameters in a structured JSON format, without requiring regular expression fallbacks or post-processing.

The most relevant candidates are summarised in Table~\ref{tab:llm_candidates}; the full list of evaluated models is provided in Appendix A. Each model was tested across two task templates (\textit{Write a Note} and \textit{Change Text Format in Word}).


\begin{table}[H]
\centering
\begin{tabular}{|l|c|c|c|c|c|c|}
\hline
\textbf{Model} & \textbf{Params} & \textbf{Quant.} & \textbf{Write} & \textbf{Format} & \textbf{Inference} & \textbf{Avg} \\
 & & & & & \textbf{Time (s)} & \textbf{Vram (\%)} \\
\hline
Gemma4 E4B         & 4B & INT4  & 100 & 100 & 21.07 & 37.53 \\
\hline
Qwen3 4B Instruct  & 4B & INT4  & 100 & 100 &  2.02 & 25.95 \\
\hline
Qwen3.5 8B         & 8B & INT4  & 100 & 100 &  8.32 & 59.67 \\
\hline
Gemma3 E2B         & 4B & INT 4 & 100 & 100 &  5.29 & 45.56 \\
\hline
Llama3.2 3B        & 3B & INT 4 & 100 & 100 &  2.33 & 17.23 \\
\hline
Ministral3 3B      & 3B & INT 4 & 100 &   0 &  2.26 & 30.90 \\
\hline
\end{tabular}
\caption{LLM candidate evaluation across two task templates.}
\label{tab:llm_candidates}
\end{table}

All six models achieved perfect intent classification on the Write a Note template, and all except Ministral3 3B generalised correctly to a more complex Format Text template. From an entirely Intention Agent perspective, the criteria to select the best model would favour the fastest and most resource-efficient model, which in this case it’s Qwen3 4B Instruct stands out with an inference time of just 2.02 s and 25.95\% average VRAM usage, while Llama3.2 3B offers similar speed at 2.33 s with the lowest VRAM footprint of all candidates (17.23\%).

However, intent classification is only one of the two roles the LLM must fulfil. The same model is also used as the Manager Agent, where it must read the full PDDL plan, maintain execution state across multiple steps, and dispatch each agent Models that performed well as Intention Agents consistently failed as Managers, they lost track of the plan, repeated steps, or didn’t follow the PDDL plan. Of all candidates evaluated, only Gemma4 E4B demonstrated reliable orchestration behaviour across all seven task templates, it’s Effective 4B architecture, gives this model a huge advantage, with far better ability to reason and manage the entire workflow.

\subsubsection{System Evaluation}
\label{sec:autopilot-system-eval}

This section presents the internal performance of the final Autopilot pipeline across all seven task templates, evaluated over 3 runs each (21 total runs). The metrics reported cover the timing breakdown of each pipeline stage, hardware resource utilisation, and the self-reported outcome of the QA Agent. Full per-run details are provided in Appendix A. For all this metrics, Gemma4 E4B to INT 4, was used and as seen in the results, was able to fully manage all workflows, with no errors.

All runs were conducted starting from a clean desktop state, with no applications open and no files pre-loaded in the target directories. This setup was chosen deliberately to maximise the number of steps evaluated per run, ensuring that the pipeline must navigate the full execution path from the initial screen state rather than benefiting from a partially prepared environment. It is important to note that all results presented in this chapter correspond to the system as evaluated at that point, without the initial screen state assessment feature described in Section~\ref{sec:system-design}; that functionality was introduced as a direct consequence of the observations derived from this comparative evaluation.

\paragraph{Manager Evalution}
\begin{table}[H]
\centering
\begin{tabular}{|l|c|c|c|c|c|c|c|}
\hline
\textbf{Template} & \textbf{Intention} & \textbf{Planner} & \textbf{Translator} & \textbf{Scout} & \textbf{Vision} & \textbf{Executor} & \textbf{QA} \\
 & \textbf{(s)} & \textbf{(s)} & \textbf{(s)} & \textbf{(s)} & \textbf{(s)} & \textbf{(s)} & \textbf{(s)} \\
\hline
Write a Note      & 11.43 &  2.61 & 1.89 &  6.36 & N/A  &  7.02 & 1.05 \\
\hline
Create Doc/Folder &  7.82 &  2.10 & 0.55 &  2.57 & 1.58 &  9.45 & 0.61 \\
\hline
Search \& Open   & 13.59 &  1.46 & 1.01 &  1.81 & 2.61 & 16.28 & 1.11 \\
\hline
Search Text       & 10.85 &  1.98 & 1.88 &  3.86 & 3.81 & 19.27 & 3.36 \\
\hline
Delete Text       & 11.94 &  1.47 & 1.69 &  8.20 & 5.01 & 27.66 & 1.08 \\
\hline
Format Text       & 15.15 &  1.38 & 3.28 &  7.43 & 4.64 & 19.25 & 3.39 \\
\hline
Copy and Paste    &  8.97 &  1.97 & 2.45 &  6.05 & 1.66 & 16.11 & 2.42 \\
\hline
Global            & 11.39 &  1.85 & 1.82 &  5.18 & 3.22 & 16.43 & 1.86 \\
\hline
\end{tabular}
\caption{Autopilot average Manager Thinking timings (seconds) per task template, averaged across 3 runs each.}
\label{tab:manager_timings}
\end{table}

Across all templates, intent classification is the main cost, averaging 11.39 s globally. This is expected, as it is the first initialization of the model and might take more time to power up and load to the VRAM. Afterwards, the model is already loaded and call the Planner Agent quickly, the same thing happens for the translator agent that is called efficiently right after the PDDL planning.

The Scout Agent time varies from nearly 2s up to almost 8s, this is due the numerous calls during more complex tasks, such as formatting a specific text in Word. Moreover, the Vision Agent is more direct, due to the straight and more natural approach. The execute action, has the highest time, this is because the Manager Agent, must think thoroughly, what the response for the Scout and Vision Agent was, and therefore give the executor the correct parameters for the given the specific situation, to execute a fallback action, a click, type or find a specific text.

Finally, the QA Agent, is always the last step of the workflow, therefore due to the explicit declaration in the prompt, the manager agent can do this step very quickly and precise.

\paragraph{Agents Evaluation}
\begin{table}[H]
\centering
\begin{tabular}{|l|c|c|c|c|c|c|c|}
\hline
\textbf{Template} & \textbf{Intention} & \textbf{Planner} & \textbf{Translator} & \textbf{Scout} & \textbf{Vision} & \textbf{Executor} & \textbf{QA} \\
 & \textbf{(s)} & \textbf{(s)} & \textbf{(s)} & \textbf{(s)} & \textbf{(s)} & \textbf{(s)} & \textbf{(s)} \\
\hline
Write a Note      & 20.41 & 0.45 &    0 & 2.39 & N/A  &  7.51 & 4.17 \\
\hline
Create Doc/Folder & 18.49 & 0.19 &    0 & 1.38 & 1.34 & 11.30 & 3.36 \\
\hline
Search \& Open   & 22.98 & 0.58 &    0 & 2.38 & 3.74 & 12.90 & 3.63 \\
\hline
Search Text       & 21.79 & 0.43 &    0 & 2.36 & 2.03 & 13.22 & 3.91 \\
\hline
Delete Text       & 25.63 & 0.47 &    0 & 2.11 & 2.04 & 15.86 & 4.58 \\
\hline
Format Text       & 31.17 & 0.45 & 0.01 & 7.43 & 4.64 & 19.66 & 4.09 \\
\hline
Copy and Paste    & 23.28 & 0.44 &    0 & 2.08 & 1.91 & 12.32 & 4.04 \\
\hline
Global            & 23.29 & 0.42 &    0 & 2.87 & 2.63 & 13.26 & 3.97 \\
\hline
\end{tabular}
\caption{Autopilot average Agent times (seconds) per task template averaged across 3 runs each.}
\label{tab:agent_timings}
\end{table}

This table breaks down the average execution time per agent across all seven templates. The Intention agent dominates the pipeline, accounting for 23.29 s globally, which represents over half of the total end-to-end time in most templates. This is expected given that it runs a quantized 4B language model locally to parse all 7 templates, from free form speech into a structured JSON. Moreover, the Intention Agent is the first step of the pipeline, therefore it might take more time to power up and load to the VRAM, which can explain the higher time compared to the rest of the agents

The next agent in line is the Planner, whichis consistently negligible at 0.42 s globally, reflecting the deterministic and lightweight nature of Fast Downward. The Translator is effectively zero across all templates, due to its pure code operation with no model inference.

Scout Agent averages 2.87 s globally, with Format Text being the highest at 7.43 s due to the greater number of visibility checks required when locating text or UI elements. Vision is N/A for Write a Note, the only template where no coordinate grounding is needed, this task only requires opening Notepad and typing. Across the remaining templates, Vision averages 2.63 s, remaining stable and predictable, due to the straightforward prompt and task it has to complete. 

Then, the Executor Agent is the second most expensive agent at 13.26 s globally, this time is static, as the system has to wait until the Windows system has updated its state, otherwise, there is a risk of capturing incorrect screenshots, or typing/clicking when the OS isn't ready yet. Finally, QA is stable across all templates at around 3.97 s globally, consistent with its fixed role and custom prompt, which is designed to be concise and direct to the point, making it easier for the model to understand and answer quickly.

\paragraph{Resources Evaluation}

\begin{table}[H]
\centering
\begin{tabular}{|l|c|c|c|c|}
\hline
\textbf{Template} & \textbf{CPU avg} & \textbf{RAM avg} & \textbf{GPU Util.} & \textbf{GPU Mem.} \\
 & \textbf{(\%)} & \textbf{(\%)} & \textbf{avg (\%)} & \textbf{avg (\%)} \\
\hline
Write a Note      & 25.3 & 70.2 & 20.9 & 81.0 \\
\hline
Create Doc/Folder & 18.6 & 78.0 & 33.0 & 80.0 \\
\hline
Search \& Open   & 25.9 & 71.4 & 23.9 & 74.1 \\
\hline
Search Text       & 25.9 & 74.1 & 22.6 & 78.8 \\
\hline
Delete Text       & 26.6 & 70.4 & 22.4 & 81.0 \\
\hline
Format Text       & 25.7 & 69.4 & 21.3 & 74.6 \\
\hline
Copy \& Paste    & 25.2 & 69.9 & 22.7 & 81.0 \\
\hline
Global            & 24.7 & 71.9 & 23.8 & 78.6 \\
\hline
\end{tabular}
\caption{Autopilot average resource usage per task template, averaged across 3 runs.}
\label{tab:resource_usage}
\end{table}

The regarding resource usage is a key aspect in this system, due to the complex organisation of the agents, it is important to reduce the resource consumption. We can see how GPU memory averages between 74\% and 81\% across templates, consistent with Ollama's managed memory, the quantized models are loaded into VRAM on demand and remain resident until unloaded, causing memory usage to stay elevated throughout execution. GPU utilisation averages between 21\% and 33\% and dropping to near zero between agent calls. CPU usage remains modest at around 25\% on average, and RAM usage stays stable between 70\% and 78\%, confirming that the system operates comfortably within consumer hardware constraints.

Finally, it is important to note, that the QA Agent reported success (YES) in 86\% of runs globally (18 out of 21). The two \textit{NO} responses occurred during the \textit{Format Text} template, where the QA Agent incorrectly judged the formatting task as incomplete despite human verification confirming success. Despite altering the prompt several times to be more concise and even telling the model to compare the final text with other non-formatted text. Meaning that the QA Agent's screenshot, despite being in high-quality, Qwen3 only process images in 1000x1000, therefore the image distortion, not only in quality but in the aspect ratio, can lead the QA Agent to be less precise.

\subsection{Comparative Evaluation}

This section presents a comparative benchmark of Autopilot against five alternative desktop automation systems. All systems were evaluated on the same seven task templates, over 3 runs each, on the same hardware described in Section 1.1. Steps across all systems are normalised to Autopilot's PDDL action schema, one step corresponds to one atomic PDDL action. This normalisation was possible because Gemini CU and Gemini CU Screenshot produce execution logs, Claude Cowork exposes a visible thought process, allowing each system's actions to be mapped to the equivalent PDDL actions and Windows Voice Access uses direct user's commands that can easily be mapped to the same schema, although there are audio recordings to verify the exact commands used.

\subsubsection{Competitor Systems}

\paragraph{Windows Voice Access}
Windows Voice Access is a built-in accessibility feature in Windows 11 that allows users to control their computer using voice commands. It provides up to 80 commands for various actions, such as opening applications, clicking buttons and typing, which are the core actions required to complete the task templates. The system was tested with the laptop's built-in microphone, and the command execution was verified through audio recordings to ensure that the correct commands were issued for each step in the task templates. Here the main disadvantage of this system is that the user is the planner and WVA is the executor, meaning that the user cannot give full prompt at the beginning of the task, but must give commands step by step, which fails to meet the requirement of a single prompt and fully autonomous execution. Moreover, the user must be familiar with the specific commands and their syntax, which can be a barrier for non-technical users.

\paragraph{Claude Cowork}
Claude Cowork is a collaborative cloud based AI system developed by Anthropic, designed to assist users in various tasks through natural language interactions. One of it's functionalities (still in Beta) is computer use. It allows the user to give a single prompt describing the task they want to accomplish, and then the system takes care of executing it, using the GUI, terminal or any available tool. The system provides a visible thought process, where it explains the steps it is taking to complete the task, which allows for a detailed evaluation of its performance and the reasoning behind its actions. 

Claude Cowork's main objective is to serve as an accessible alternative to Claude Code for non-technical users, providing a dedicated Windows application with a more intuitive interface. However, during computer use sessions, the system uses other tools, such as opening the terminal or calling the Windows Run dialog (Win+R), actions that despite being effective, they may generate distrust in users who are not familiar.

Unlike Autopilot's PDDL architecture, Claude Cowork operates reactively, meaning that it captures a screenshot, decides an action, executes it, and repeats the cycle until the goal is completed, with no first planning stage. This approach is great for simple tasks, which are very straightforward and don't require any long and complex reasoning, however, as evidenced in the recorded runs and later reflected in the comparative metrics, it also leads to non-optimal execution paths in more complex tasks, where the system becomes repetitive after reaching dead ends or fails to identify the most direct route to the objective. 

A recurring pattern observed across templates is that Claude Cowork opens a target file or window and immediately minimises it, as documented in Run 3 of the \textit{Search \& Open} template, meaning that in the next step, Claude thinks the task was not completed, despite being completed. Although, the system's fast actions allow it to recover and complete most tasks within a competitive time.

\paragraph{OpenClaw}
OpenClaw is an open-source desktop automation system designed for the model to use various tools at hand to comply with user demands, without providing a dedicated GUI. It was evaluated using Gemma4 E4B INT4 and Qwen3.5 8B INT4 across two tool integration approaches: skills and MCPs, tested independently.

Two skills were evaluated: Desktop Control [ref], which covers mouse, keyboard, screen capture, and clipboard management via OpenCV, and Windows UI Automation [ref], which automates Windows GUI interactions through native APIs and PowerShell. Two MCPs were also tested: PyAutoGUI MCP [ref], which exposes mouse, keyboard, and screen control as model-callable tools, and Windows-MCP by CursorTouch [ref], which combines PyAutoGUI with the Windows UIAutomation framework for native element interaction without computer vision. None of these configurations was sufficient for the model to reliably plan and execute the task templates.

\paragraph{Gemini Computer Use}
Gemini Computer Use (Gemini CU) is a cloud-based desktop automation system developed by Google, that allows users to control the desktop through natural language commands. However, it currently only offers an API with no dedicated user interface, which makes it less accessible for non-technical users. Thus, a Python script had to be developed in order to interact with the model. Both variants used \textit{gemini-3-flash-preview} as the underlying model.

Two variants were implemented and evaluated, due to the specific recommendation in the Google documentation, which states it is better for the developer to provide a screenshot after each action, rather than leaving the model to decide when to capture the screen. Although, this goes against the core principle of fully autonomous execution, therefore, both approaches were implemented and evaluated to determine the best integration pattern for this system.

In the first variant, Gemini CU, \textit{take\_screenshot} is exposed as one of the available tools, the model decides autonomously when to observe the screen, calling it only when it deems necessary. In the second variant, Gemini Screenshot (Gemini SS), the screenshot tool is removed entirely and the orchestration script automatically captures and injects a screenshot into the conversation after every action, ensuring the model always has an up-to-date view of the desktop state before its next decision.

A significant practical limitation observed throughout both evaluations was API availability. Out of 42 total runs, 22 had to be repeated due to the API returning resource exhaustion errors, indicating that the service was under heavy load. This instability introduces a substantial reliability concern for any deployment that depends on consistent response times.


\subsubsection{Evaluation Metrics}

Eleven metrics were defined across four categories: performance, reliability, variability, and privacy. Table~\ref{tab:metrics} summarises each metric, its description, and its possible values or range.

\begin{table}[H]
\centering
\small
\begin{tabular}{|l|p{3.5cm}|p{5cm}|p{3cm}|}
\hline
\textbf{ID} & \textbf{Name} & \textbf{Description} & \textbf{Values / Range} \\
\hline
M-01 & Task Success Rate (TSR) & Average proportion of runs completed successfully & 0 to 1 \\
\hline
M-02 & Step Accuracy & Fraction of executed steps that were error-free & 0 to 1 \\
\hline
M-03 & End-to-End Latency & Average time from prompt to completion, if time exceed 5 minutes the task is aborted. & seconds \\
\hline
M-04 & Planning Error Rate & Fraction of steps with wrong or missing actions & 0 to 1 \\
\hline
M-05 & Tool Hallucination Rate & Fraction of steps invoking a non-existent or incorrect tool & 0 to 1 \\
\hline
M-06 & Visual Hallucination Rate & Fraction of steps where the vision model misidentified on-screen elements & 0 to 1 \\
\hline
M-07 & Total Error Rate & Combined error rate (sum of M-04, M-05, M-06) & 0 to 1 \\
\hline
M-08 & Standard Deviation of Steps & Variability of step count across repetitions, the higher the deviation, the less consistent & Float \\
\hline
M-09 & Standard Deviation of Latency & Variability of latency across repetitions, the higher the deviation the less consistent & Float \\
\hline
M-10 & Data Exposure Level & Amount of sensitive data sent outside the device & 
      \begin{tabular}[c]{@{}l@{}}
        0: No data leaves the device \\
        1: Metadata or task text sent externally \\
        2: Sensitive data sent externally \\
        3: Sensitive data used to train models \\
      \end{tabular} \\
\hline
M-11 & Total Steps & Total number of steps taken to complete a run & integer \\
\hline
\end{tabular}
\caption{Evaluation metrics used in the comparative benchmark.}
\label{tab:metrics}
\end{table}

\subsubsection{Full System Comparison}

Table~\ref{tab:full_comparative} presents all six systems evaluated across the same seven task templates with three runs each, providing an unconstrained baseline before any filtering is applied.

\begin{table}[H]
\centering
\begin{tabular}{|p{4.5cm}|c|c|c|c|c|c|}
\hline
\textbf{Metric} & \textbf{Autopilot} & \textbf{WVA} & \textbf{OpenClaw} & \textbf{Claude} & \textbf{Gemini} & \textbf{Gemini} \\
 & & & & \textbf{CW} & \textbf{CU} & \textbf{SS} \\
\hline
M-01 Task Success Rate (TSR)        & 1.000  & 1.000 & 0.000 & 1.000 & 0.333 & 0.952 \\
\hline
M-02 Step Accuracy                  & 0.925  & 0.968 & 0.000 & 0.775 & 0.479 & 0.893 \\
\hline
M-03 End-to-End Latency             & 111.83 & 47.00 & - & 111.75 & 47.90 & 77.67 \\
\hline
M-04 Planning Error Rate            & 0.000  & 0.000 & 0.000 & 0.225 & 0.358 & 0.107 \\
\hline
M-05 Tool Hallucination Rate        & 0.000  & 0.000 & 0.000 & 0.000 & 0.163 & 0.000 \\
\hline
M-06 Visual Hallucination Rate      & 0.075  & 0.032 & 0.000 & 0.000 & 0.000 & 0.000 \\
\hline
M-07 Total Error Rate               & 0.075  & 0.032 & 0.000 & 0.225 & 0.521 & 0.107 \\
\hline
M-08 Standard Deviation of Steps    & 0.000  & 0.034 & 0.000 & 0.139 & 0.588 & 0.191 \\
\hline
M-09 Standard Deviation of Latency  & 0.123  & 0.095 & 0.000 & 0.152 & 0.556 & 0.527 \\
\hline
M-10 Data Exposure Level            & 0      & 0     & 0     & 3     & 3     & 3     \\
\hline
M-11 Total Steps                    & 4.57   & 10.90 & 0.000 & 6.38  & 8.48  & 5.86  \\
\hline
\end{tabular}
\caption{Full comparative evaluation across all six systems with three runs per template.}
\label{tab:full_comparative}
\end{table}

A cross-cutting observation that emerged from reviewing all competitor systems is that none of them performs an initial assessment of the current desktop state before committing to a plan or action. Every system evaluated, whether reactive like Claude CW and Gemini CU or command-driven like WVA, assumes a blank slate and begins issuing actions immediately, without first checking what is already visible on screen. This means they cannot take advantage of elements that are already open, nor detect early that the required conditions are not yet met. After analysing these results, it was decided to extend Autopilot with an initial screen state check: before invoking the Planning Agent, the Manager Agent now calls the Scout Agent to assess what is currently visible on the desktop, allowing the pipeline to ground its plan in the actual context from the very first step. This feature was not present during the experiments reported in this chapter and is therefore not reflected in the metrics above; it represents a direct improvement derived from the comparative analysis.

This section focuses on WVA and Gemini SS, the two systems excluded from the next comparison, and how they relate to Autopilot. The remaining systems (OpenClaw, Claude CW, and Gemini CU) are discussed in the following sections.

WVA achieves a perfect TSR (1.000) and actually outperforms Autopilot on step accuracy (0.968 vs 0.925) and total error rate (0.032 vs 0.075), with a lower latency of 47.00 s. However, these numbers come at the cost of autonomy: WVA requires the user to issue commands step by step, which is why it needs 10.90 steps on average, more than double Autopilot's 4.57. It is not a fair autonomous competitor, but it serves as a bottom line of current alternatives to assistive technologies. Moreover, WVA has as default settings to mantain everything locally with no data being sent across servers although this feature can be turned on at any time.

Gemini SS scores TSR 0.952 and a total error rate of 0.107, making it the second most accurate system. The constant scrennshot feedback injected after each action clearly helps the model keep track of its progress and future actions, resulting in lower planning errors (0.107) and no tool or visual hallucinations. Despite this, Gemini SS is excluded from the later autonomous comparison because the screenshot injection is managed externally by the orchestration script, not by the model itself, which breaks the core requirement of fully autonomous execution. Additionally, Gemini SS sends information to external servers, where they use this data to better train the models.

\subsubsection{Autonomous Systems}

Table~\ref{tab:autonomous_comparative} restricts the comparison to systems that run fully autonomously, excluding WVA and Gemini SS for reason mentioned before. Therefore, we can directly compare the core autonomous capabilities of Autopilot, OpenClaw, Claude CW, and Gemini CU.

\begin{table}[H]
\centering
\begin{tabular}{|p{4.5cm}|c|c|c|c|}
\hline
\textbf{Metric} & \textbf{Autopilot} & \textbf{OpenClaw} & \textbf{Claude} & \textbf{Gemini} \\
 & & & \textbf{CW} & \textbf{CU} \\
\hline
M-01 Task Success Rate (TSR)        & 1.000  & 0.000 & 1.000 & 0.333 \\
\hline
M-02 Step Accuracy                  & 0.925  & 0.000 & 0.775 & 0.479 \\
\hline
M-03 End-to-End Latency             & 111.83 & - & 111.75 & 47.90 \\
\hline
M-04 Planning Error Rate            & 0.000  & 0.000 & 0.225 & 0.358 \\
\hline
M-05 Tool Hallucination Rate        & 0.000  & 0.000 & 0.000 & 0.163 \\
\hline
M-06 Visual Hallucination Rate      & 0.075  & 0.000 & 0.000 & 0.000 \\
\hline
M-07 Total Error Rate               & 0.075  & 0.000 & 0.225 & 0.521 \\
\hline
M-08 Standard Deviation of Steps    & 0.000  & 0.000 & 0.139 & 0.588 \\
\hline
M-09 Standard Deviation of Latency  & 0.123  & 0.000 & 0.152 & 0.556 \\
\hline
M-10 Data Exposure Level            & 0      & 0     & 3     & 3     \\
\hline
M-11 Total Steps                    & 4.57   & 0.000 & 6.38  & 8.48  \\
\hline
\end{tabular}
\caption{Autonomous systems comparison: Autopilot, OpenClaw, Claude CW, and Gemini CU.}
\label{tab:autonomous_comparative}
\end{table}

This section focuses on Claude CW and Gemini CU, the two systems excluded from the final comparison, and how they relate to Autopilot. OpenClaw is also present but fails entirely (TSR 0.000), and is discussed in depth in the following section alongside Autopilot.

Claude CW matches Autopilot on TSR (1.000) but accumulates a planning error rate of 0.225, compared to Autopilot's zero, which is consistent with its reactive architecture: no upfront plan means the model occasionally chooses the wrong action or gets stuck. This also translates into more steps on average (6.38 vs 4.57) and a higher step standard deviation (0.139 vs 0.000). On the positive side, Claude CW produces no tool or visual hallucinations, meaning its errors are purely reasoning-based rather than perceptual. However, as a cloud system, it transmits screenshots and conversation context on every action (M-10 = 3), which rules it out for privacy-sensitive deployments.

Gemini CU performs significantly worse, achieving only TSR 0.333 and accumulating errors across all three categories: planning (0.358), tool hallucination (0.163), and a total error rate of 0.521. It is also the least consistent system, with a step standard deviation of 0.588 and a latency standard deviation of 0.556. Like Claude CW, it carries a data exposure level of 3. Both cloud systems are excluded from the final comparison for these reasons. The next section isolates Autopilot and OpenClaw as the only two fully local candidates.

\subsubsection{Fully Local Systems}

Table~\ref{tab:local_comparative} isolates the two systems that run entirely on local hardware without any cloud dependency, enabling a direct comparison of on-device automation performance.

\begin{table}[H]
\centering
\begin{tabular}{|p{4.5cm}|c|c|}
\hline
\textbf{Metric} & \textbf{Autopilot} & \textbf{OpenClaw} \\
\hline
M-01 Task Success Rate (TSR)        & 1.000  & 0.000 \\
\hline
M-02 Step Accuracy                  & 0.925  & 0.000 \\
\hline
M-03 End-to-End Latency             & 111.83 & - \\
\hline
M-04 Planning Error Rate            & 0.000  & 0.000 \\
\hline
M-05 Tool Hallucination Rate        & 0.000  & 0.000 \\
\hline
M-06 Visual Hallucination Rate      & 0.075  & 0.000 \\
\hline
M-07 Total Error Rate               & 0.075  & 0.000 \\
\hline
M-08 Standard Deviation of Steps    & 0.000  & 0.000 \\
\hline
M-09 Standard Deviation of Latency  & 0.123  & 0.000 \\
\hline
M-10 Data Exposure Level            & 0      & 0     \\
\hline
M-11 Total Steps                    & 4.57   & 0.000 \\
\hline
\end{tabular}
\caption{Fully local systems comparison: Autopilot and OpenClaw.}
\label{tab:local_comparative}
\end{table}

The fully local comparison yields the starkest result in the entire benchmark: Autopilot completes every task (TSR 1.000) while OpenClaw completes none (TSR 0.000). Both systems run on the same hardware, use comparable quantised models, and share a data exposure score of zero, so the gap cannot be attributed to hardware, cloud access, or privacy trade-offs.

The difference is purely architectural. Autopilot separates intention parsing, symbolic planning, visual grounding, and execution into specialised agents, each with a focused role and a tailored prompt. OpenClaw relies on a single general-purpose LLM to reason about the full task, decide which tool to invoke, and recover from failures, all in one unstructured loop. None of the tool configurations tested, whether skills or MCPs, were sufficient to overcome this limitation. The results show that access to the right tools is not enough without a planning layer that structures how they are used.

This directly supports one of the central claims of this project: running local models is necessary for privacy and offline operation, but it is not sufficient for reliable automation. The structured pipeline is the critical component, and Autopilot demonstrates that strong performance on consumer hardware is achievable without any cloud dependency.

\section{Qualitative Evaluation}

In addition to the quantitative benchmark, a qualitative user study was conducted to assess real-world perceived usability and accessibility. Ten participants were invited to interact with Autopilot on the same hardware used for the comparative evaluation described in Section~\ref{sec:autopilot-system-eval}. Each participant issued a free-form voice command of their choice, covering six of the seven supported task templates, and then completed a Google Forms survey of six Likert-scale questions (1 = worst, 5 = best) alongside an optional free-text comment. All ten participants confirmed that the system completed their requested task successfully.

Table~\ref{tab:qualitative_responses} reports the individual responses of each participant. Column headers are abbreviated; their full descriptions are: \textit{Recovery} (ability to recover from errors), \textit{Behaviour} (absence of detours or unexpected actions), \textit{Privacy} (perceived importance of local processing), \textit{Accessibility} (suitability for users with motor difficulties), \textit{Daily Use} (likelihood of daily use), and \textit{vs Others} (likelihood of choosing Autopilot over other alternatives).

\begin{table}[H]
\centering
\begin{tabular}{|c|l|c|c|c|c|c|c|}
\hline
\textbf{\#} & \textbf{Task} & \textbf{Recovery} & \textbf{Behaviour} & \textbf{Privacy} & \textbf{Accessibility} & \textbf{Daily Use} & \textbf{vs Others} \\
\hline
1  & Open File      & 5 & 2 & 4 & 4 & 4 & 4 \\
\hline
2  & Format Text    & 4 & 2 & 5 & 5 & 3 & 3 \\
\hline
3  & Format Text    & 3 & 4 & 2 & 2 & 2 & 4 \\
\hline
4  & Open File      & 5 & 5 & 5 & 5 & 4 & 4 \\
\hline
5  & Copy \& Paste  & 5 & 4 & 2 & 4 & 4 & 3 \\
\hline
6  & Create Folder  & 4 & 5 & 5 & 4 & 4 & 3 \\
\hline
7  & Write Note     & 5 & 5 & 4 & 4 & 4 & 4 \\
\hline
8  & Search Text    & 4 & 4 & 4 & 4 & 3 & 4 \\
\hline
9  & Write Note     & 5 & 4 & 4 & 5 & 5 & 4 \\
\hline
10 & Copy \& Paste  & 5 & 5 & 5 & 4 & 5 & 3 \\
\hline
\textbf{Mean} & & \textbf{4.5} & \textbf{4.0} & \textbf{4.0} & \textbf{4.1} & \textbf{3.8} & \textbf{3.6} \\
\hline
\end{tabular}
\caption{Qualitative user study responses. All ratings on a 1--5 scale (1 = worst, 5 = best).}
\label{tab:qualitative_responses}
\end{table}

Error recovery received the highest mean score (4.5), indicating that participants consistently perceived the system as capable of handling mistakes without failing the task. One participant noted: \textit{Very impressive how despite identifying wrong certain elements in the screen it detects that and is able to continue to next step.} Behavioural predictability averaged 4.0, the two lower scores (both 2) corresponded to participants who attempted the Format Text template, where a higher number of intermediate steps and the visual grounding behaviour may have appeared less direct to a first-time user. Privacy scored 4.0 on average, though opinions were notably different, with participants of a more technical background rating it highest, while one noted \textit{``it is not very important to me the data in my device''}, reflecting that the perceived value of local processing depends on the user. Accessibility suitability for motor impaired people averaged 4.1, with one outlier (score 2) corresponding to the same participant who reported transcription errors.

Likelihood of daily use (3.8) and preference over alternatives (3.6) were the lowest scoring dimensions, consistent with the system's current limitation to seven task templates and the latency due to the complex orquestration in local INT4 inference. Two participants specifically mentioned execution time, with one commenting \textit{``It took a while, it wasn't quick.''} Speech recognition accuracy was the most frequently cited area for improvement, two participants noted that the system did not transcribe their command correctly on the first attempt, although in both cases the task was completed successfully.
